\documentclass[10pt,a4paper]{amsart}
\usepackage[margin=19mm]{geometry}
\usepackage{amsmath,amssymb,mathtools}
\usepackage[scr=rsfs,cal=euler]{mathalfa}
\usepackage{xfrac}
\usepackage[unicode,hidelinks,bookmarks=false]{hyperref}
\newcommand{\ie}{i.e.\ }
\newcommand{\cf}{cf.\ }
\newcommand{\resp}{respectively\ }
\newcommand{\Subsection}{\subsection}
\providecommand{\nobreakdash}{\nobreak\hbox{-}}
\setlength{\emergencystretch}{2em}
\multlinegap0pt
\usepackage{bm}
\usepackage{bbm}
\usepackage{dsfont}
\usepackage{xspace}
\usepackage{cancel}
\usepackage{mathtools}
\usepackage{amsthm}
\makeatletter
\@ifundefined{lemma}{\newtheorem{lemma}{Lemma}}{}
\@ifundefined{theorem}{\newtheorem{theorem}{Theorem}}{}
\makeatother
\title[The completeness and congruences of quasi-Boolean algebras]{The completeness and congruences of quasi-Boolean algebras}
\author{Xiaohao Liu, Heyan Wang,, Wenjuan Chen}
\date{Source: arXiv:2510.23094v1 (CC BY 4.0)}
\begin{document}
\maketitle

\noindent\emph{Source and licence.} The completeness and congruences of quasi-Boolean algebras. Version arXiv:2510.23094v1 (CC BY 4.0), distributed under the Creative Commons Attribution 4.0 International licence, as recorded in the arXiv licence metadata for this exact version and shown on the abstract page https://arxiv.org/abs/2510.23094v1. Full rights notice: licenses/arxiv-2510.23094-CC-BY-4.0.txt.\par\smallskip

\noindent\textit{Editorial scope.} Counted unit: the Theorem whose source label is \texttt{Proposition 4.1} in the article (source0.tex lines 479--484, source label \texttt{Proposition 4.1}), delivered together with the article's own complete original proof of it (source0.tex lines 486--515, one \texttt{proof} environment of 30 lines). No mathematics is written, repaired or re-derived: statement and proof are the article's own text line for line. Required context retained uncounted: Lemma 3 (lines 445--447). Every definition, display and label of those lines is retained unchanged. Omitted from this package and not redistributed: 1--444, 448--478, 485--485, 516--1073. Nothing inside a retained excerpt is omitted or altered: every retained body is delivered exactly as the article's lines are written, so no omission receipt of any kind accompanies this package. Citations: the retained excerpts carry no citation command at all, so no bibliography entry of the article is transcribed and no reference is redistributed. Reference adaptation: none was needed and none was made; every reference inside the delivered unit resolves inside it, so no unresolved reference or citation is produced. Printed slips kept literal: no printed word, symbol, hypothesis or number of the counted statement or of its proof was altered. Layout and class: the article's own class is \texttt{amsart} with packages including amsrefs, mathdots, blkarray, multicol, graphicx, caption, amscd, upgreek, stmaryrd, longtable, fontenc, latexsym, amsmath, amssymb; the wrapper is the lane's pinned \texttt{amsart} editorial wrapper at 10pt on A4, which restates the theorem-like environments the retained text needs and adds the article's own macro definitions that the retained text needs (none), with extra wrapper packages (bm, bbm, dsfont, xspace, cancel, mathtools). The delivered numbering is the unit's own. Assets: no retained span contains a figure environment, a graphics-inclusion command, a PDF-page inclusion, a table or a TikZ picture; no image file is copied into this package, the package declares no asset dependency and no image file is redistributed with it. Licence: version arXiv:2510.23094v1 of this article is distributed under the Creative Commons Attribution 4.0 International licence, as recorded in the arXiv licence metadata for this exact version and shown on the abstract page https://arxiv.org/abs/2510.23094v1. A full rights notice is delivered at licenses/arxiv-2510.23094-CC-BY-4.0.txt. Characters: every retained excerpt is pure ASCII, so the delivered engine, XeLaTeX with its default Latin Modern text font, reports no missing character.
\begin{lemma}\label{Lemma 3}\label{vee 0}
Let $\mathbf{Q}$ be a flat QB-algebra and $x,y\in Q$. Then $x\wedge y=x\vee y=0$.
\end{lemma}

\begin{theorem}\label{Proposition 4.1}
Let $p$ and $q$ be terms in the language of QB-algebras. Then
$$
\mathbb{FQB}\models p\approx q  \; \mathrm{iff}  \; \mathrm{\mathbf{F}_{3}}\models p\approx q.
$$
\end{theorem}

\begin{proof}
Suppose that $\mathbb{FQB}\models p\approx q$. Then we have $\mathrm{\mathbf{F}_{3}}\models p\approx q$.

Conversely, suppose that $\mathrm{\mathbf{F}_{3}}\models p\approx q$ and if $\mathbb{FQB}\nvDash p\approx q$, then there exists a flat QB-algebra $\mathbf{F}$ such that $\mathrm{\mathbf{F}}\nvDash p\approx q$. We distinguish several cases.

(1) Both $p(x_{1},x_{2},\cdots,x_{n})$ and $q(x_{1},x_{2},\cdots,x_{m})$ contain at least an occurrence of a binary operation. Then for any $a_{1},a_{2},\cdots,a_{m},\cdots,a_{n}\in F$, we have that
$$
p^\mathbf{F}(a_{1},a_{2},\cdots,a_{n})=0^\mathbf{F}=q^\mathbf{F}(a_{1},a_{2},\cdots,a_{m})
$$
from Lemma \ref{Lemma 3}, against the hypothesis.

(2) Either $p$ or $q$ contains at least an occurrence of a binary operation. Without loss of generality, we assume that $p$ contains at least an occurrence of a binary operation and $q$ does not contain any occurrence of a binary operation. Then $q$ is one of the following forms.

(2.1) If $q$ is a variable $\emph{x}$ and contains $k\,(0\leq k)$ occurrences of the operation $^{*}$, then we can assign $x$ the value $c\in \mathbf{F}_{3}$ and the variables $x_{1},\cdots ,x_{n}$ any values $0,c,d\in \mathbf{F}_{3}$, then $p^\mathbf{F}_{3}(0,\cdots,c,\cdots,d)=0$ $\neq $ $q^\mathbf{F}_{3}(c)$, this is a contradiction with $\mathrm{\mathbf{F}_{3}}\models p\approx q$.

(2.2) If $q$ is a constant and contains $k\,(0\leq k)$ occurrences of the operation $^{*}$, then $q^\mathbf{F}=0$, it follows that $q^\mathbf{F}=0=p^\mathbf{F}$, against the hypothesis.

(3) Neither $p$ nor $q$ contains any occurrence of a binary operation, then $p$ and $q$ contain at most one variable, it turns out that they have one of the following forms.

(3.1) $p$ is a constant and contains $k \,(0\leq k)$ occurrences of the operation $^{*}$, $q$ is a constant and contains $l \, (0\leq l)$ occurrences of the operation $^{*}$. Then we have $p^\mathbf{F}=0=q^\mathbf{F}$, against the hypothesis.

(3.2) $p$ is a variable $\emph{x}$ and contains $k\,(0\leq k)$ occurrences of the operation $^{*}$, $q$ is a constant and contains $l\, (0\leq l)$ occurrences of the operation $^{*}$. Assign $\emph{x}$ the value $c\in \mathbf{F}_{3}$, we can get that $p^\mathbf{F}_{3}(c)\neq 0=q^\mathbf{F}_{3}$, this is a contradiction with
$\mathrm{\mathbf{F}_{3}}\models p\approx q$.

(3.3) $p$ is a variable $x$ and contains $k\,(0\leq k)$ occurrences of the operation $^{*}$, $q$ is a variable $y$ and contains $l\,(0\leq l)$ occurrences of the operation $^{*}$. Assign $x$ the value $c\in \mathbf{F}_{3}$ and $y$ the value $0\in \mathbf{F}_{3}$, it follows that $p^\mathbf{F}_{3}(c)\neq q\mathbf{F}_{3}(0)$, this is a contradiction with $\mathrm{\mathbf{F}_{3}}\models p\approx q$.

(3.4) $\emph{p}$ is a variable $\emph{x}$ and contains $k\,(0\leq k)$ occurrences of the operation $^{*}$, $\emph{q}$ is a variable $\emph{x}$ and contains $l\,(0\leq l)$ occurrences of the operation $^{*}$. If $\emph{k}$ and $\emph{l}$ have same parity, then $p^\mathbf{F}$=$q^\mathbf{F}$, against the hypothesis. If $\emph{k}$ and $\emph{l}$ have different parities, then we can assign $\emph{x}$ the value $c\in \mathbf{F}_{3}$, it follows that $p^\mathbf{F}_{3}(c)\neq q^\mathbf{F}_{3}(c)$, this is a contradiction with $\mathrm{\mathbf{F}_{3}}\models p\approx q$.

Thus if $\mathrm{\mathbf{F}_{3}}\models p\approx q$, then we have $\mathbb{FQB}\models p\approx q$.
\end{proof}

\end{document}
